\section{\bt{Use-Time Enforcement}{使用时的强制保证}}
\label{sec:execution}
\subsection{\bt{Revision checks}{修订核对}}
\bi{Agents declare the metric keys and revisions used by a query; the middleware checks those references at execution and rejects absent, revoked, inaccessible, or superseded revisions. SQL review checks the supported grain and join requirements. Explicit rechecking at use makes a definition's maintained state actionable, and revocation notices reach agents on their next tool response. Canonical compilation supplies the calculation encoded by the recorded fields.}
{智能体声明查询所用的指标键和修订号，中间层在执行时核对这些引用，拒绝不存在、已撤销、不可访问或已被替代的修订。SQL 审查检查支持的粒度与连接要求。使用时的显式重查使定义的维护状态直接约束执行，撤销通知通过智能体下一次工具响应传达。规范编译提供记录字段所表达的计算。}

\subsection{\bt{Snapshot-bound execution}{绑定快照的执行}}
\label{sec:snapshot}
\bi{Revision checks ensure that a declared revision was valid when it was last maintained, not that its conditions hold on the data the query reads. Between the two, a writer may commit, and the default version source may not yet show the write. Snapshot-bound execution closes this gap. For a \code{run\_sql} call that declares revisions, \system\ (1) opens a repeatable-read, read-only transaction and reads the transactional versions; (2) collects the conditions of the declared revisions; (3) for each condition, looks up a verdict at its identity and the snapshot's versions of its tables, and otherwise runs $Q_c$ inside the transaction; (4) rejects the call if any condition fails, which also hands the definition to regular maintenance; and (5) otherwise executes the business SQL in the same transaction.}
{修订核对只能保证声明的修订在最近一次维护时有效，不能保证其条件在查询读取的数据上成立。两者之间可能有写入提交，而默认版本来源可能还没反映这次写入。绑定快照的执行填补了这一空隙。对声明了修订的 \code{run\_sql} 调用，\system{}（1）开启一个可重复读的只读事务并读取事务性版本；（2）收集所声明修订的条件；（3）对每个条件，按其身份及快照中其依赖表的版本查找结论，没有就在该事务内执行 $Q_c$；（4）任一条件不成立即拒绝调用，并把定义交给常规维护；（5）否则在同一事务中执行业务 SQL。}

\bi{Reuse in step~(3) is sound because PostgreSQL snapshots are nested: a snapshot taken later sees every committed transaction that an earlier one saw. Every write to a table increments its transactional version inside the writing transaction, so two snapshots that read equal versions of a table see the same committed writes to it, and a verdict computed on one holds on the other. Verdicts computed by maintenance outside a snapshot qualify only if the versions read immediately before and after the check are equal, so that no write committed while the check ran. Concurrent identical checks at the same versions run once and are shared. The cost is a counter update per writing statement, which Section~\ref{sec:enforce-eval} measures.}
{第（3）步的复用是可靠的，因为 PostgreSQL 的快照是嵌套的：后取的快照能看到先取快照看到的每个已提交事务。对一张表的每次写入都在写事务内增加其事务性版本，因此读到该表相同版本的两个快照看到的对它的已提交写入完全相同，在一个快照上算出的结论在另一个快照上同样成立。维护在快照之外算出的结论，只有在检查前后各读一次的版本相同、即检查期间没有写入提交时，才可用于快照。相同版本上的相同检查并发到达时只执行一次并共享结果。代价是每条写语句多一次计数更新，第~\ref{sec:enforce-eval} 节测量这一代价。}

\subsection{\bt{Scope of the guarantee}{保证的范围}}
\bi{The guarantee covers declared uses: a query that declares a revision runs only on data on which the revision's conditions hold. It does not make undeclared SQL follow a definition, and it does not certify that an agent's own SQL matches the canonical calculation; SQL review checks the supported grain and join requirements, and canonical compilation supplies the calculation for agents that use it. Section~\ref{sec:scenarios} reports how often agents declare definitions and follow them.}
{这一保证覆盖声明了修订的使用：声明某个修订的查询只在该修订条件成立的数据上执行。它不能使未声明的 SQL 遵循定义，也不能证明智能体自己写的 SQL 与规范计算一致；SQL 审查检查支持的粒度与连接要求，规范编译为使用它的智能体提供计算。第~\ref{sec:scenarios} 节报告智能体声明并遵循定义的频率。}

